\documentclass[10pt,a4paper]{amsart}
\usepackage[margin=19mm]{geometry}
\usepackage{amsmath,amssymb,mathtools}
\usepackage[scr=rsfs,cal=euler]{mathalfa}
\usepackage{xfrac}
\usepackage[unicode,hidelinks,bookmarks=false]{hyperref}
\newcommand{\ie}{i.e.\ }
\newcommand{\cf}{cf.\ }
\newcommand{\resp}{respectively\ }
\newcommand{\Subsection}{\subsection}
\providecommand{\nobreakdash}{\nobreak\hbox{-}}
\setlength{\emergencystretch}{2em}
\multlinegap0pt
\usepackage{amssymb}
\usepackage{bm}
\usepackage{enumerate}
\newtheorem{lemma}{Lemma}
\newcommand{\M}{\mathcal{M}}
\newcommand{\crit}{\operatorname{crit}}
\newcommand{\dn}{\operatorname{D}}
\newcommand{\up}{\operatorname{U}}
\title{$2$-colourability of the maximum ranked elements of a combinatorially sphere-like ranked poset}
\author{Anupam Mondal, Sajal Mukherjee, Pritam Chandra Pramanik}
\date{Source: arXiv:2604.08210v1 (CC BY 4.0)}
\begin{document}
\maketitle

\noindent\emph{Source and licence.} $2$-colourability of the maximum ranked elements of a combinatorially sphere-like ranked poset. Version arXiv:2604.08210v1 (CC BY 4.0), distributed by arXiv under the Creative Commons Attribution 4.0 International licence, http://creativecommons.org/licenses/by/4.0/, as recorded in the arXiv licence metadata for this exact version and shown on the abstract page https://arxiv.org/abs/2604.08210v1. Full licence notice: \texttt{licenses/arxiv-2604.08210-CC-BY-4.0.txt}.\par\smallskip

\noindent\textit{Editorial scope.} Counted unit: the article's lemma bd (source0.tex lines 203--205). The article's complete original proof of it is delivered with it (source0.tex lines 206--250, one \texttt{proof} environment of 45 lines). No mathematics is written, repaired or re-derived: the statement and the proof are the article's own lines, in the article's own notation, and no retained line is substituted or re-flowed.  Retained uncounted as the source-local context the proof needs: lines 140--142; lines 159--161; lines 176--179.  Nothing was omitted from any retained span, so the package carries no omission record.  No retained line contains a citation, so the wrapper carries no bibliography and no bibliography entry.  No retained line contains a figure environment, an image-inclusion command or drawing code, so no image is delivered and the package declares no asset dependency.  Editorial formatting only, in addition: an AMS \texttt{amsart} preamble that restates the article's own declarations and macros used by the retained lines, namely the environments \texttt{lemma} on separate counters, together with the standard packages those lines need.  The retained environments print in this document as Lemma 1 in order of appearance, whereas the article prints the same items with the numbering of its own full document.  The complete native source of the article as distributed in the arXiv e-print for this exact version is bundled unchanged as \texttt{sources/198141/source0.tex}.
	\begin{lemma}\label{poincare}
		Let  $S$ be sphere-like poset of rank $k$, then $d_{k-1}\circ d_{k}=0$.
	\end{lemma}

	 \begin{lemma}\label{ft}
	 	Let $S$ be a ranked poset and $\M$ be any acyclic matching on it. Then, for any $x\in \crit_q^{(\M)}$ and $z\in \up_{q-1}^{(\M)}$, we have  $\sum_{y\in \nabla(z)} \ell^{(\M)}_{(x,y)}=0$.
	 \end{lemma}

     \begin{lemma}\label{dn}
     	Let $S$ be a sphere-like poset of rank $k$, and $\M$ be any acyclic matching on
     	 $S$. Let $ \dn_{k-1}^{(\M)}=\{x_1,\ldots,x_n\}$. For any $\sigma=\sum_{i=1}^n t_ix_i ~(\in C_{k-1})$, if $d_{k-1}(\sigma)=0$, then $\sigma=0$.
     \end{lemma}

     \begin{lemma} \label{bd}
     	Let $S$ be a sphere-like ranked poset of rank $k$, then for any $\eta\in C_{k-1}$, such that $d_{k-1}(\eta)=0$, there exists $\xi\in C_k$ such that $d_k(\xi)=\eta$.
     \end{lemma}

     \begin{proof}
     	Let $\M$ be an acyclic matching on $S$, such that there exists no $\M$-critical element of rank $(k-1)$. Let $\eta=\sum_{i=1}^{n} s_i x_i$, where $\up_{k-1}^{(\M)}=\{x_1,\ldots,x_m\}$, $\dn_{k-1}^{(\M)}=\{x_{m+1},\ldots,x_{n}\}$, and $s_i\in \mathbb{Z}_2$, for all $i \in \{1,\ldots,n\}$.  Let $(x_j,y_j)\in \M$, for all $j\in \{1,\ldots,m\}$.	 Note that, for all $j\in\{1,\ldots,m\}$, $\M_j=\M\setminus\{(x_j,y_j)\}$ is also an acyclic matching on $S$.
     	For any $j\in \{1,\ldots,m\}$, we define $\overrightarrow{y_j}^{(\M_j)}$ ($\in C_k$), as follows.
     	\[ \overrightarrow{y_j}^{(\M_j)}:= \sum_{y\in S_k} \ell^{(\M_j)}_{(y_j,y)}\cdot y.\]
     
     	We show that, for $\xi= \sum_{j=1}^m s_j\overrightarrow{y_j}^{(\M_j)}$, $d_k(\xi)=\eta$. For any $j\in\{1,\ldots,m\}$, we have
     	\begin{align} \label{c}
     		\nonumber	d_k (\overrightarrow{y_j}^{(\M_j)}) =&d_k(\sum_{y\in S_{k}}\ell^{(\M_j)}_{(y_j,y)}\cdot y)\\
     		\nonumber	= & \sum_{y\in S_{k}}\ell^{(\M_j)}_{(y_j,y)}\cdot d_k(y)\\
     		\nonumber	= &\sum_{y\in S_{k}}\ell^{(\M_j)}_{(y_j,y)}\cdot \sum_{i=1}^n \chi_{(y,x_i)} \cdot x_i\\
     		\nonumber	= & \sum_{i=1}^n\sum_{y\in S_{k}}\ell^{(\M_j)}_{(y_j,y)}\cdot \chi_{(y,x_i)} \cdot x_i\\
     	    \nonumber		= & \sum_{i=1}^m\sum_{y\in S_{k}}\ell^{(\M_j)}_{(y_j,y)}\cdot \chi_{(y,x_i)} \cdot x_i+ \sum_{i=m+1}^n\sum_{y\in S_{k}}\ell^{(\M_j)}_{(y_j,y)}\cdot \chi_{(y,x_i)} \cdot x_i\\
     	    		=& \sum_{i=1}^m \sum_{y\in \nabla(x_i)}\ell^{(\M_j)}_{(y_j,y)} \cdot x_i	+ \sum_{i=m+1}^n \sum_{y\in \nabla(x_i)}\ell^{(\M_j)}_{(y_j,y)} \cdot x_i.	
     	\end{align}
   Note that, since $x_j\lessdot y_j$, there exists only one trajectory $P$, which is the trivial trajectory $P: y_j$, such that $\operatorname{init}(P)=y_j$ and $\operatorname{term}(P)\in \nabla(x_j)$, otherwise it violates the acyclicity of $\M$. In this case, $\operatorname{init}(P)=\operatorname{term}(P)=y_j$. Hence, $\sum_{y\in \nabla(x_j)}\ell^{(\M_j)}_{(y_j,y)}=1$. Now, since $y_j \in \crit_k^{(\M_j)}$, and, for any $x_i$, such that $i\in \{1,\ldots,m\}\setminus \{j\}$,  $x_i\in \up_{k-1}^{(\M_j)}$, from Lemma~\ref{ft}, we have $\sum_{y\in \nabla(x_i)}\ell^{(\M_j)}_{(y_j,y)}=0$. Hence Equation~\eqref{c} becomes
     	
     	\begin{align}  \label{a} 
     		d_k (\overrightarrow{y_j}^{(\M_j)})	= x_j +   \sum_{i=m+1}^n \sum_{y\in \nabla(x_i)}\ell^{(\M_j)}_{(y_j,y)} \cdot x_i	.
     	\end{align}
 
  Now,		
     	\begin{align}	\label{b}
     	\nonumber   d_k ( \sum_{j=1}^m s_j \overrightarrow{y_j}^{(\M_j)})=&\sum_{j=1}^m s_jd_k (\overrightarrow{y_j}^{(\M_j)}) \\
     		\nonumber	=&	\sum_{j=1}^m s_j \left(x_j +   \sum_{i=m+1}^n \sum_{y\in \nabla(x_i)}\ell^{(\M_j)}_{(y_j,y)} \cdot x_i	 \right) \text{ (from Equation~\eqref{a})}\\
     		\nonumber	=&	\sum_{j=1}^m s_j x_j +  \sum_{j=1}^m s_j \left( \sum_{i=m+1}^n \sum_{y\in \nabla(x_i)}\ell^{(\M_j)}_{(y_j,y)} \cdot x_i \right)\\
     		\nonumber	=&	\sum_{j=1}^m s_j x_j +   \sum_{i=m+1}^n \left( \sum_{j=1}^m  \sum_{y\in \nabla(x_i)}s_j \cdot \ell^{(\M_j)}_{(y_j,y)}  \right)\cdot x_i\\
     		\nonumber	=&	\sum_{i=1}^m s_i x_i +   \sum_{i=m+1}^n \left( \sum_{j=1}^m  \sum_{y\in \nabla(x_i)}s_j \cdot \ell^{(\M_j)}_{(y_j,y)}  \right)\cdot x_i\\
     		\nonumber	=&	\sum_{i=1}^m s_i x_i + \sum_{i=m+1}^n s_ix_i+  \sum_{i=m+1}^n \left( \sum_{j=1}^m  \sum_{y\in \nabla(x_i)}s_j \cdot \ell^{(\M_j)}_{(y_j,y)} +s_i \right)\cdot x_i\\
     		\nonumber	=&	\eta +  \sum_{i=m+1}^n t_i \cdot x_i \text{ (where, $t_i=\sum_{j=1}^m  \sum_{y\in \nabla(x_i)}s_j \cdot \ell^{(\M_j)}_{(y_j,y)} +s_i$)}\\
     		\implies  d_k ( \sum_{j=1}^m s_j \overrightarrow{y_j}^{(\M_j)})+\eta=& \sum_{i=m+1}^n t_i \cdot x_i.
     		\end{align}
     		
     		Now, since 
     		\begin{align*}
     		d_{k-1}(\sum_{i=m+1}^n t_i \cdot x_i)=&d_{k-1}( d_k ( \sum_{j=1}^m s_j \overrightarrow{y_j}^{(\M_j)})+\eta)\\
     		=& 	d_{k-1}( d_k ( \sum_{j=1}^m s_j \overrightarrow{y_j}^{(\M_j)}))+d_{k-1}(\eta)\\
     		=&	d_{k-1}( d_k ( \sum_{j=1}^m s_j \overrightarrow{y_j}^{(\M_j)}))\\
     		=& 0 \text{ (by Lemma~\ref{poincare})}, 
     		\end{align*}
     by Lemma~\ref{dn}, $\sum_{i=m+1}^n t_i \cdot x_i=0$. Hence from Equation~\eqref{b}, we get
    \begin{align*}
    	  &d_k ( \sum_{j=1}^m s_j \overrightarrow{y_j}^{(\M_j)})+\eta=0\\
    	  \implies &d_k ( \sum_{j=1}^m s_j \overrightarrow{y_j}^{(\M_j)})=\eta.
    \end{align*}	
     \end{proof}

\end{document}
